\paragraph{El mito de los dos diodos y la barrera invisible}
Luis dibuja en la pizarra el esquema mental que acaba de armar: un transistor NPN conectado a una batería de \qty{12}{\volt} entre colector (positivo) y emisor (negativo). Observa que el terminal positivo roba electrones del colector, dejando iones positivos cerca de la unión. 

``Si la base es P y el colector es N, y el colector está a \qty{12}{\volt}, esa unión está en polarización inversa. Es un muro infranqueable'', razona Luis. ``Y si el emisor de alguna forma empuja electrones a la base, estos simplemente caerían en los huecos. ¡Es como tener dos diodos opuestos! No debería fluir corriente hacia el colector''.

Carl, escuchando atentamente, aplaude. ``Tu lógica es perfecta para componentes discretos, Luis. Si sueldas dos diodos por sus ánodos, jamás tendrás un transistor. El secreto del transistor no reside solo en el dopaje, sino en un diseño físico extremo''.

\paragraph{Las dos zonas de deplexión}
Carl toma una tiza y dibuja el interior del transistor. Le señala a Luis que, al haber tres bloques, existen \emph{dos} zonas de deplexión: la unión Base-Emisor y la unión Base-Colector. 


Con los \qty{12}{\volt} en el colector y el emisor a tierra (y la base desconectada), la unión Base-Colector sufre una polarización inversa masiva. Su zona de deplexión se ensancha enormemente. El colector se vuelve un imán gigante esperando cargas negativas, pero la barrera no deja que nada cruze.

\paragraph{Los \qty{0.7}{\volt} y la Base microscópica}
``¿Qué pasa si ahora le aplicamos \qty{0.7}{\volt} a la base respecto al emisor?'', pregunta Carl. 

``Esa es la barrera de potencial del silicio'', responde Luis casi por instinto. ``La unión Base-Emisor se polariza en directa. Su zona de deplexión colapsa y el emisor, que es una \emph{metralleta} fuertemente dopada, empieza a disparar millones de electrones libres hacia la base''.

``¡Exacto! Y aquí es donde el diseño rompe las reglas'', sonríe Carl. ``La base tiene huecos, sí, pero el fabricante la hace \emph{extremadamente delgada} (apenas unos micrómetros) y la dopa \emph{muy débilmente}. Cuando esa avalancha de electrones entra a la base, hay tan pocos huecos disponibles que solo el 1\% de los electrones logra chocar, recombinarse y salir por el cable de la base''.

\paragraph{El salto al vacío}
Carl dibuja flechas masivas cruzando la base de largo. ``El 99\% restante de los electrones entra con tanta energía cinética que cruzan la minúscula base sin encontrar un solo hueco. Antes de que puedan darse cuenta, se asoman al abismo: la enorme zona de deplexión de la unión Base-Colector''.

Luis abre los ojos sorprendido al ver el panorama completo. ``¡Claro! Esa zona estaba en inversa para bloquear los huecos de la base... pero para un electrón libre que acaba de invadir la zona a toda velocidad, ¡es un tobogán electromagnético!''.

``¡Bingo!'', exclama Carl. ``Esa tensión de \qty{12}{\volt} crea un campo eléctrico intenso. Cuando los electrones huérfanos del emisor sobreviven al cruce de la base y tocan el borde de esa zona de deplexión, el campo eléctrico del colector los atrapa y los barre violentamente hacia el terminal positivo de la batería''.
